%% kapa: exam=1 round=27 year=2016 period=1 form=A subject=economics subject_ko=경제학원론 seq=36 number=76
%% source: Q-Net 감정평가사 기출문제 (cst00309) artlSeq=5169402 file=extracted/1008/5169402/제27회 감정평가사 1차 문제/제27회 감정평가사 1교시문제(A형).pdf
%% conversion: kordoc 4.24.51 -> tools/gen1.py (auto); review=proofread-sonnet; flags=none
\begin{question}{76}
통화량(M)을 현금(C)과 요구불예금(D)의 합으로, 본원통화(B)를 현금(C)과 지급준비금(R)의 합으로 정의하자. 이 경우 현금보유비율($cr$)은 C/D, 지급준비금 비율($rr$)은 R/D로 나타낼 수 있다. 중앙은행이 본원통화를 공급할 때 민간은 현금 보유분을 제외하고는 모두 은행에 예금하며, 은행은 수취한 예금 중 지급준비금을 제외하고는 모두 대출한다고 가정한다. $cr$이 0.2, $rr$이 0.1이면 통화승수의 크기는?
\begin{choices}
  \item 1.5
  \item 2.0
  \item 3.7
  \item 4.0
  \item 5.3
\end{choices}
\end{question}
\answer{4}
